\documentclass{article}
\usepackage[T1]{fontenc}
\usepackage{lmodern}
\usepackage{graphicx}
\usepackage{amsmath,amssymb}
\usepackage[a4paper,margin=2cm]{geometry}
\usepackage[absolute,overlay]{textpos}
\setlength{\TPHorizModule}{1mm}
\setlength{\TPVertModule}{1mm}
\usepackage[indonesian]{babel}
\usepackage{hyperref}
\setlength{\emergencystretch}{2em}
% unit: Halaman 23

\begin{document}

% === Halaman 23 (python/native) ===

Contoh 3:   a  b = (x3 + x2 + 1)  (x2 + x )  = x5 + 2x4 + x3 + x2 + x (mod 2)





                 = x5 + x3 + x2 + x = 10110

 Karena m = 4 hasilnya direduksi menjadi derajat < 4 oleh  sebuah irreducible 

polynomial  f(x) = x4 + x + 1 


    Jadi,   (x5 + x3 + x2 + x)  mod f(x)   =  x3  = 1000

  a  b = 1000 


Rinaldi Munir/II4021 Kriptografi

23

Proses pembagiannya ditunjukkan sebagai berikut:


                      x

           x4 + x + 1       x5 + x3 + x2 + x 



       x5 +        x2 + x   +    



               x3      →  sisa pembagian

\end{document}
